\documentclass[11pt]{article}

\usepackage[margin=1in]{geometry}
\usepackage[T1]{fontenc}
\usepackage[utf8]{inputenc}
\usepackage{lmodern}
\usepackage{microtype}
\usepackage{amsmath,amssymb,mathtools}
\usepackage{hyperref}
\usepackage{enumitem}
\usepackage{parskip}

\hypersetup{
  colorlinks=true,
  linkcolor=blue,
  urlcolor=blue,
  pdftitle={proto.py: Theoretical and Mathematical Overview}
}

\setlist[itemize]{nosep,leftmargin=2em}
\setlist[enumerate]{nosep,leftmargin=2em}

\newcommand{\R}{\mathbb{R}}

\title{\texttt{proto.py}: Theoretical and Mathematical Overview}
\date{}

\begin{document}
\maketitle

This note explains \texttt{proto.py} as a model, not as a software artifact. It avoids a code-path walkthrough and instead describes the mathematical objects, training objectives, and generative factorization implemented by the file.

\section{High-level picture}

The file implements a two-stage generative model:

\begin{enumerate}
  \item A sparse-latent autoencoder learns to map an image into a structured latent representation built from dictionary atoms.
  \item An autoregressive transformer learns a prior over that sparse latent representation.
\end{enumerate}

The basic stage-1 pipeline is
\[
x \xrightarrow{E_\theta} z_e \xrightarrow{\text{sparse bottleneck}} z_q \xrightarrow{G_\psi} \hat x.
\]

For stage 2, the learned sparse codes are flattened into a sequence:
\[
\text{stage-1 sparse codes} \xrightarrow{\text{flatten}} y_{1:T}
\xrightarrow{\text{Transformer prior}} p(y_{1:T}).
\]

The central idea is that the latent space is not modeled as an unconstrained dense tensor. Instead, each latent site or latent patch is constrained to lie on a sparse dictionary manifold.

\section{Notation}

Let:
\begin{itemize}
  \item $x \in [-1,1]^{3 \times R \times R}$ be an input image;
  \item $E_\theta$ be the encoder;
  \item $G_\psi$ be the decoder;
  \item $z_e = E_\theta(x) \in \R^{C \times H \times W}$ be the encoder latent;
  \item $z_q$ be the sparse-structured latent produced by the bottleneck;
  \item $K$ be the number of dictionary atoms;
  \item $s$ be the sparsity level;
  \item $D \in \R^{C \times K}$ be the learned atom dictionary in the per-location case.
\end{itemize}

For the patch-based case, the dictionary lives in a higher-dimensional patch space:
\[
D_p \in \R^{(C p^2) \times K}.
\]

The latent sequence modeled by the transformer has length
\[
T = H W D_t,
\]
where $D_t$ is the token depth per spatial location:
\begin{itemize}
  \item $D_t = s$ in the real-valued coefficient regime;
  \item $D_t = 2s$ in the quantized coefficient regime.
\end{itemize}

\section{Stage 1: sparse-latent autoencoding}

\subsection{Encoder-decoder backbone}

The encoder-decoder pair is a multiscale convolutional autoencoder with residual blocks and optional attention. Conceptually, it defines two maps:
\[
E_\theta : \R^{3 \times R \times R} \to \R^{C \times H \times W},
\]
\[
G_\psi : \R^{C \times H \times W} \to \R^{3 \times R \times R}.
\]

In the default $128 \times 128$ configuration used here, the target latent grid is $8 \times 8$, so the encoder must reduce spatial resolution by a factor of $16$ along each axis. If $L$ is the number of spatial downsampling steps, then
\[
H = W = \frac{R}{2^L}.
\]
The implementation uses $L = 4$, so with $R = 128$,
\[
H = W = \frac{128}{2^4} = 8.
\]

\subsubsection{Downsampling law}

Each encoder downsampling layer pads on the right and bottom by one pixel and then applies a $3 \times 3$ convolution with stride $2$. If the incoming spatial size is $n$ and $n$ is even, then the output size is
\[
n'
=
\left\lfloor \frac{(n+1)-3}{2} \right\rfloor + 1
=
\frac{n}{2}.
\]

So the spatial path is exactly
\[
128 \to 64 \to 32 \to 16 \to 8.
\]

The decoder mirrors this with nearest-neighbor interpolation by a factor of $2$, followed by a $3 \times 3$ convolution, so its spatial path is
\[
8 \to 16 \to 32 \to 64 \to 128.
\]

\subsubsection{Channel schedule in this implementation}

The code does not use an unrestricted doubling schedule. Instead it defines
\[
\mathrm{ch\_mult}_\ell = \min(2^\ell, 2), \qquad \ell = 0,\dots,L,
\]
so with base width $\mathrm{num\_hiddens} = 128$ and $L = 4$, the channel counts are
\[
128 \cdot (1,2,2,2,2) = (128,256,256,256,256).
\]

Thus the encoder expands channels once, from $128$ to $256$, and then keeps the pyramid at $256$ channels until the final projection to the bottleneck width $\mathrm{embedding\_dim} = 16$.

\subsubsection{Residual block as the basic nonlinear operator}

At each resolution, the main computation is performed by ResNet blocks. If $h \in \R^{C_{\mathrm{in}} \times r \times r}$, one block has the form
\[
\mathrm{ResBlock}(h)
=
S(h)
+
W_2 \,
\mathrm{Drop}\!\left(
\phi\!\left(
N_2\!\left(
W_1 \phi(N_1(h))
\right)
\right)
\right),
\]
where $N_1, N_2$ are group-normalization operators, $\phi$ is the SiLU nonlinearity, $W_1, W_2$ are $3 \times 3$ convolutions, and $S$ is either the identity or a learned $1 \times 1$ / $3 \times 3$ shortcut if the channel dimension changes.

\subsubsection{Attention at the bottleneck scale}

The model also includes self-attention blocks. For a feature map $h \in \R^{C \times r \times r}$, let
\[
Q \in \R^{r^2 \times C},
\qquad
K \in \R^{C \times r^2},
\qquad
V \in \R^{C \times r^2}
\]
be $1 \times 1$ convolutions of the normalized feature map after reshaping the spatial grid into a sequence of length $r^2$. Define the attention matrix
\[
A = \mathrm{softmax}\!\left(\frac{QK}{\sqrt{C}}\right) \in \R^{r^2 \times r^2}.
\]

Then the attention block is
\[
\mathrm{Attn}(h)
=
h
+
W_o (V A^\top).
\]

Two details matter: attention can optionally be inserted at user-chosen resolutions via \texttt{attn\_resolutions}, and there is always a middle attention block at the lowest-resolution feature map. So even without high-resolution attention, the model still gets one global mixing step at the $8 \times 8$ bottleneck scale.

\subsubsection{Exact encoder geometry for $128 \times 128 \to 8 \times 8$}

With \texttt{num\_residual\_layers = 2}, \texttt{num\_hiddens = 128}, \texttt{num\_downsamples = 4}, and \texttt{embedding\_dim = 16}, the encoder shape progression is:
\begin{itemize}
  \item input image: $x \in \R^{3 \times 128 \times 128}$;
  \item input $3 \times 3$ convolution: $\R^{3 \times 128 \times 128} \to \R^{128 \times 128 \times 128}$;
  \item level 0 residual stack: stays at $\R^{128 \times 128 \times 128}$;
  \item downsample: $\R^{128 \times 128 \times 128} \to \R^{128 \times 64 \times 64}$;
  \item level 1 residual stack: $\R^{128 \times 64 \times 64} \to \R^{256 \times 64 \times 64}$;
  \item downsample: $\R^{256 \times 64 \times 64} \to \R^{256 \times 32 \times 32}$;
  \item level 2 residual stack: stays at $\R^{256 \times 32 \times 32}$;
  \item downsample: $\R^{256 \times 32 \times 32} \to \R^{256 \times 16 \times 16}$;
  \item level 3 residual stack: stays at $\R^{256 \times 16 \times 16}$;
  \item downsample: $\R^{256 \times 16 \times 16} \to \R^{256 \times 8 \times 8}$;
  \item level 4 residual stack plus middle $\mathrm{Res}$-$\mathrm{Attn}$-$\mathrm{Res}$ block: stays at $\R^{256 \times 8 \times 8}$;
  \item output normalization, SiLU, and final $3 \times 3$ projection: $\R^{256 \times 8 \times 8} \to \R^{16 \times 8 \times 8}$.
\end{itemize}

So the encoder output is
\[
z_e \in \R^{16 \times 8 \times 8}.
\]

This is a spatial compression by $16$ along each axis and a dense coordinate compression from
\[
3 \cdot 128 \cdot 128 = 49152
\]
input scalars to
\[
16 \cdot 8 \cdot 8 = 1024
\]
latent scalars, i.e. a $48:1$ reduction in raw coordinate count before the sparse dictionary bottleneck is applied.

\subsubsection{Exact decoder geometry for $8 \times 8 \to 128 \times 128$}

The decoder is not a U-Net: it has no encoder-decoder skip connections. All information that reaches image space must pass through the latent tensor $z_q \in \R^{16 \times 8 \times 8}$.

The decoder first lifts the $16$ latent channels back to $256$ channels with a $3 \times 3$ convolution, applies the same middle $\mathrm{Res}$-$\mathrm{Attn}$-$\mathrm{Res}$ computation at $8 \times 8$, and then runs a multiscale upsampling tower. Because the decoder uses $\mathrm{num\_res\_blocks} + 1$ residual blocks per scale, it is slightly heavier than the encoder on the way back up.

Its shape progression is:
\begin{itemize}
  \item latent input: $z_q \in \R^{16 \times 8 \times 8}$;
  \item input projection: $\R^{16 \times 8 \times 8} \to \R^{256 \times 8 \times 8}$;
  \item middle $\mathrm{Res}$-$\mathrm{Attn}$-$\mathrm{Res}$: stays at $\R^{256 \times 8 \times 8}$;
  \item up level at $8 \times 8$: stays at $\R^{256 \times 8 \times 8}$, then upsample to $\R^{256 \times 16 \times 16}$;
  \item up level at $16 \times 16$: stays at $\R^{256 \times 16 \times 16}$, then upsample to $\R^{256 \times 32 \times 32}$;
  \item up level at $32 \times 32$: stays at $\R^{256 \times 32 \times 32}$, then upsample to $\R^{256 \times 64 \times 64}$;
  \item up level at $64 \times 64$: stays at $\R^{256 \times 64 \times 64}$, then upsample to $\R^{256 \times 128 \times 128}$;
  \item final level at $128 \times 128$: maps $\R^{256 \times 128 \times 128} \to \R^{128 \times 128 \times 128}$;
  \item output normalization, SiLU, and final $3 \times 3$ RGB projection: $\R^{128 \times 128 \times 128} \to \R^{3 \times 128 \times 128}$.
\end{itemize}

If \texttt{out\_tanh=True}, the final image is
\[
\hat x = \tanh(G_\psi(z_q)) \in [-1,1]^{3 \times 128 \times 128}.
\]

So mathematically the decoder is a learned map from an $8 \times 8$ latent field back to image space, with all high-resolution detail reconstructed from the compressed representation rather than copied through skips.

The key constraint is what happens between encoder and decoder: the dense latent $z_e$ is projected onto a sparse dictionary model before decoding.

\subsection{Per-location sparse bottleneck}

In the non-patch formulation, each spatial latent vector
\[
z_e(u) \in \R^C, \qquad u \in \{1,\dots,H\} \times \{1,\dots,W\}
\]
is represented as a sparse linear combination of atoms:
\[
z_q(u) = \sum_{j=1}^{s} a_j(u)\, d_{i_j(u)},
\]
where $d_k \in \R^C$ is the $k$-th column of $D$, $i_j(u) \in \{1,\dots,K\}$ are selected atom indices, and $a_j(u) \in \R$ are sparse coefficients.

Equivalently, if $\alpha(u) \in \R^K$ is a sparse coefficient vector, then
\[
z_q(u) = D \alpha(u), \qquad \|\alpha(u)\|_0 \le s.
\]

The sparse coding problem at each site is
\[
\alpha^*(u)
=
\arg\min_{\alpha \in \R^K}
\|z_e(u) - D\alpha\|_2^2
\quad
\text{subject to}
\quad
\|\alpha\|_0 \le s.
\]

\subsection{Orthogonal Matching Pursuit}

The sparse approximation is computed with batched Orthogonal Matching Pursuit (OMP). Given a signal $z \in \R^C$, initialize
\[
r^{(0)} = z, \qquad S^{(0)} = \emptyset.
\]

At iteration $m = 1,\dots,s$:
\begin{enumerate}
  \item Select the atom with largest correlation to the current residual:
  \[
  k^{(m)} = \arg\max_{k \notin S^{(m-1)}} |d_k^\top r^{(m-1)}|.
  \]
  \item Expand the support:
  \[
  S^{(m)} = S^{(m-1)} \cup \{k^{(m)}\}.
  \]
  \item Refit the coefficients on the active support:
  \[
  \alpha_{S^{(m)}}^{(m)}
  =
  \arg\min_{\beta}
  \|z - D_{S^{(m)}} \beta\|_2^2.
  \]
  \item Update the residual:
  \[
  r^{(m)} = z - D_{S^{(m)}} \alpha_{S^{(m)}}^{(m)}.
  \]
\end{enumerate}

After $s$ steps, OMP returns a support of size exactly $s$ and aligned coefficients. Two details matter conceptually: the support size is fixed, and the selected slots are reordered by descending coefficient magnitude before stage 2 sees them.

\subsection{Dictionary geometry}

The dictionary atoms are constrained to unit norm:
\[
\|d_k\|_2 = 1 \quad \text{for all } k.
\]

This separates direction from magnitude: atom identity carries the latent direction, while the coefficient carries the scale and sign. The implementation also monitors dictionary coherence,
\[
\mu(D) = \max_{i \ne j} |d_i^\top d_j|,
\]
which measures the largest absolute off-diagonal cosine similarity.

\subsection{Straight-through sparse projection}

The sparse coding map is non-differentiable in practice because OMP involves hard support selection. The model therefore uses a straight-through surrogate. Let $z_q$ be the sparse reconstruction obtained from OMP. The latent passed to the decoder is
\[
z_{\text{ste}} = z_e + \mathrm{sg}(z_q - z_e),
\]
where $\mathrm{sg}$ denotes stop-gradient.

Forward pass:
\[
z_{\text{ste}} = z_q.
\]

Backward pass:
\[
\frac{\partial z_{\text{ste}}}{\partial z_e} = I.
\]

So the decoder sees the sparse-projected latent, while the encoder receives an identity-style backward signal through the bottleneck.

\subsection{Stage-1 objective}

The stage-1 loss has two parts: image reconstruction and bottleneck alignment.

The image reconstruction term is
\[
\mathcal{L}_{\text{recon}}
=
\|\hat x - x\|_2^2,
\qquad
\hat x = G_\psi(z_{\text{ste}}).
\]

The bottleneck term has a dictionary-side part and an encoder-side commitment part:
\[
\mathcal{L}_{\text{dict}}
=
\|z_q - \mathrm{sg}(z_e)\|_2^2,
\]
\[
\mathcal{L}_{\text{commit}}
=
\|\mathrm{sg}(z_q) - z_e\|_2^2.
\]

The bottleneck loss is
\[
\mathcal{L}_{\text{bottleneck}}
=
\mathcal{L}_{\text{dict}}
+
\beta \mathcal{L}_{\text{commit}},
\]
where $\beta$ is the commitment weight. The full stage-1 objective is
\[
\mathcal{L}_{\text{stage1}}
=
\mathcal{L}_{\text{recon}}
+
\lambda_b \mathcal{L}_{\text{bottleneck}},
\]
where $\lambda_b$ is the external bottleneck weight.

\subsection{What is and is not differentiated}

The sparse coefficients and supports are obtained by solving a hard sparse-coding problem. Training does not differentiate through the combinatorial support-selection map itself. Instead:
\begin{itemize}
  \item the encoder is optimized through the straight-through approximation;
  \item the dictionary is optimized through the reconstruction induced by the selected support;
  \item support selection is treated as frozen inside each forward pass.
\end{itemize}

\section{Tokenization of sparse codes}

Stage 2 does not see dense latents. It sees a tokenized version of the sparse representation. At each latent site $u$, the sparse code consists of
\[
\{(i_1(u), a_1(u)), \dots, (i_s(u), a_s(u))\}.
\]

\subsection{Real-valued coefficient regime}

In the real-valued regime, only atom identities are discretized into tokens:
\[
t_j(u) = i_j(u), \qquad j=1,\dots,s.
\]

The coefficients remain continuous side information:
\[
c_j(u) = a_j(u).
\]

So each spatial location emits $s$ discrete tokens and $s$ real-valued scalars, and the stage-2 sequence length becomes
\[
T = H W s.
\]

The vocabulary contains $K$ atom tokens, one padding token, and one beginning-of-sequence token.

\subsection{Quantized coefficient regime}

In the quantized regime, each coefficient is discretized into one of $B$ bins:
\[
q(a_j(u)) \in \{1,\dots,B\}.
\]

The token stream alternates atom identity and coefficient bin:
\[
(i_1(u), q(a_1(u)), i_2(u), q(a_2(u)), \dots, i_s(u), q(a_s(u))).
\]

So each location emits $2s$ tokens and the sequence length becomes
\[
T = H W (2s).
\]

The vocabulary contains $K$ atom tokens, $B$ coefficient-bin tokens, one padding token, and one beginning-of-sequence token.

\subsection{Coefficient clipping}

When coefficients are used as decoder inputs in the real-valued regime, they are projected into a bounded interval:
\[
\Pi_{[-c_{\max}, c_{\max}]}(a) = \min(\max(a,-c_{\max}), c_{\max}).
\]

This defines the actual bounded sparse-code manifold used by the decoder.

\section{Coefficient quantization}

\subsection{Uniform quantization}

Let $a \in [-c_{\max}, c_{\max}]$. The uniform map is
\[
u(a) = \frac{a + c_{\max}}{2c_{\max}} \in [0,1],
\]
followed by discretization into $B$ bins:
\[
b(a) = \mathrm{round}\left(u(a)(B-1)\right).
\]

Decoding maps the bin index back to a fixed bin center.

\subsection{Mu-law quantization}

To allocate more resolution near zero, the model also supports mu-law companding. Define the normalized coefficient
\[
\bar a = a / c_{\max} \in [-1,1].
\]

Then apply
\[
f_\mu(\bar a)
=
\mathrm{sign}(\bar a)
\frac{\log(1 + \mu |\bar a|)}{\log(1+\mu)}.
\]

This companded value is quantized uniformly in $[-1,1]$, then inverted at decode time.

\section{Patch-based sparse bottleneck}

Here the sparse object is not a single latent vector $z_e(u)$, but a local latent patch.

\subsection{Patch extraction}

Let
\[
P_u(z_e) \in \R^{C p^2}
\]
be the flattened latent patch centered around location $u$, with patch size $p$ and stride $r$. The sparse coding problem becomes
\[
\alpha^*(u)
=
\arg\min_{\alpha}
\|P_u(z_e) - D_p \alpha\|_2^2
\quad
\text{subject to}
\quad
\|\alpha\|_0 \le s.
\]

The reconstructed patch is
\[
\hat P_u = D_p \alpha^*(u).
\]

\subsection{Why patches change the representation}

Patch-based coding treats a small local neighborhood as the primitive object, allowing a single atom to represent a structured motif in latent space rather than one isolated feature vector. That changes both the geometry and the sequence semantics.

\subsection{Patch reconstruction operators}

Two overlap rules are implemented.

\paragraph{Center-crop stitching.}
Each reconstructed patch contributes only its central $r \times r$ region. If the patch size is $p$ and stride is $r$, define the crop margin
\[
c = \frac{p-r}{2}.
\]

Each reconstructed patch $\hat P_u$ is reshaped into a $C \times p \times p$ tensor, cropped to
\[
\hat P_u[:, c:c+r, c:c+r],
\]
and tiled into the output grid without averaging.

\paragraph{Hann overlap-add.}
Let $w \in \R^{p \times p}$ be a separable Hann window. Each reconstructed patch is weighted by $w$ and overlap-added:
\[
\tilde z = \sum_u W_u \odot \hat P_u,
\]
where $W_u$ is the shifted window at patch location $u$. The final reconstruction divides by the accumulated weight map:
\[
z_q = \frac{\sum_u W_u \odot \hat P_u}{\sum_u W_u + \varepsilon}.
\]

This is a weighted partition-of-unity style reconstruction and behaves especially well at 50\% overlap.

\section{Stage 2: autoregressive prior over sparse codes}

Once stage 1 is trained, it defines a sparse-code distribution over the training set. Stage 2 fits an autoregressive prior to that induced representation.

\subsection{Sequence geometry}

The sparse code tensor has shape
\[
H \times W \times D_t,
\]
where $D_t$ is the token depth per spatial position. Flattening gives a sequence
\[
y_{1:T}, \qquad T = H W D_t.
\]

The ordering is raster order over spatial sites with sparse-slot depth varying within each site.

\subsection{Transformer factorization}

In the purely discrete case, the model is
\[
p(y_{1:T}) = \prod_{t=1}^{T} p(y_t \mid y_{<t}).
\]

In the real-valued coefficient case, the stage-2 representation consists of atoms and coefficients:
\[
(a_1,c_1), \dots, (a_T,c_T).
\]

The model factorizes this as
\[
p(a_{1:T}, c_{1:T})
=
\prod_{t=1}^{T}
p(a_t \mid a_{<t}, c_{<t})
\,
p(c_t \mid a_t, a_{<t}, c_{<t}).
\]

\subsection{Positional structure in the transformer}

Each sequence position corresponds to a spatial site index, a depth-slot index, and a token type. If position $t$ corresponds to spatial index $s(t)$ and depth slot $d(t)$, then the input embedding is conceptually
\[
e_t
=
e_{\text{token}}(y_t)
+
e_{\text{space}}(s(t))
+
e_{\text{depth}}(d(t))
+
e_{\text{type}}(t).
\]

\section{Stage 2 in the quantized regime}

\subsection{Likelihood}

Let $y_t$ be the next token. The transformer produces logits $\ell_t \in \R^V$ over the vocabulary. The model distribution is
\[
p(y_t = v \mid y_{<t})
=
\frac{\exp(\ell_{t,v})}{\sum_{v'} \exp(\ell_{t,v'})}.
\]

\subsection{Loss}

Training minimizes the negative log-likelihood:
\[
\mathcal{L}_{\text{CE}}
=
-
\sum_{t=1}^{T}
\log p(y_t \mid y_{<t}).
\]

\subsection{Structural masking}

Because the quantized representation alternates atom tokens and coefficient-bin tokens, valid outputs depend on the parity of the position inside each local sparse stack. So the support of $p(y_t \mid y_{<t})$ depends on whether $t$ is an atom step or a coefficient step.

\section{Stage 2 in the real-valued regime}

\subsection{Autoregressive conditioning}

At training time, the transformer receives a shifted atom sequence
\[
[\text{BOS}, a_1, a_2, \dots, a_{T-1}]
\]
and a shifted coefficient sequence
\[
[0, c_1, c_2, \dots, c_{T-1}].
\]

The hidden state at time $t$ therefore summarizes
\[
(a_{<t}, c_{<t}).
\]

From that state the model predicts a categorical distribution for $a_t$ and a scalar coefficient for $c_t$ conditioned on the chosen atom.

\subsection{Atom-conditioned coefficient normalization}

For atom $k$, compute empirical statistics $(\mu_k, \sigma_k)$ from stage-1 sparse codes. The normalized coefficient target is
\[
\tilde c_t
=
\mathrm{clip}
\left(
\frac{c_t - \mu_{a_t}}{\sigma_{a_t}},
-M,
M
\right),
\]
where $M$ is a fixed normalization bound.

The model predicts $\hat{\tilde c}_t$, and decoding maps it back via
\[
\hat c_t
=
\hat{\tilde c}_t \sigma_{a_t} + \mu_{a_t}.
\]

\subsection{Coefficient head factorization}

Let $h_t$ be the transformer hidden state before the output heads. The coefficient predictor uses both the historical context summarized by $h_t$ and the embedding of the selected atom $a_t$:
\[
\hat{\tilde c}_t = f_{\text{coeff}}(h_t, e(a_t)).
\]

This matches the factorization
\[
p(c_t \mid a_t, a_{<t}, c_{<t}).
\]

\section{Stage-2 objectives in the real-valued regime}

The real-valued path always includes atom cross-entropy and then adds one auxiliary coefficient objective.

\subsection{Atom cross-entropy}

The categorical term is
\[
\mathcal{L}_{\text{atom}}
=
-
\sum_{t=1}^{T} \log p(a_t \mid a_{<t}, c_{<t}).
\]

\subsection{Direct coefficient regression}

Two direct scalar objectives are supported.

Mean squared error:
\[
\mathcal{L}_{\text{coeff-MSE}}
=
\sum_{t=1}^{T} (\hat{\tilde c}_t - \tilde c_t)^2.
\]

Huber loss:
\[
\mathcal{L}_{\text{coeff-Huber}}
=
\sum_{t=1}^{T}
\mathrm{Huber}_\delta(\hat{\tilde c}_t - \tilde c_t).
\]

\subsection{Reconstruction-aware coefficient losses}

Let
\[
R(a_{1:T}, c_{1:T})
\]
be the operator that reshapes a sequence into the latent sparse grid and reconstructs the corresponding latent tensor through the stage-1 bottleneck. The target sparse latent is
\[
z^* = R(a_{1:T}, c_{1:T}).
\]

Predicted atoms plus predicted coefficients:
\[
\hat z = R(\hat a_{1:T}, \hat c_{1:T}).
\]
\[
\mathcal{L}_{\text{recon-MSE}}
=
\|\hat z - z^*\|_2^2.
\]

Ground-truth atoms plus predicted coefficients:
\[
\hat z = R(a_{1:T}, \hat c_{1:T}),
\]
\[
\mathcal{L}_{\text{gt-atom-recon-MSE}}
=
\|\hat z - z^*\|_2^2.
\]

\subsection{Full real-valued stage-2 loss}

The full loss is
\[
\mathcal{L}_{\text{stage2}}
=
\mathcal{L}_{\text{atom}}
+
\lambda_c \mathcal{L}_{\text{coeff}},
\]
where $\mathcal{L}_{\text{coeff}}$ is one of the coefficient objectives above.

\section{Scheduled sampling}

Teacher forcing trains on the true history
\[
(a_{<t}, c_{<t}),
\]
but generation uses model-sampled history. If $\rho(\tau)$ is the replacement probability at training progress $\tau$, then for a previous time step $m < t$ the training input is
\[
(\tilde a_m, \tilde c_m)
=
\begin{cases}
(\hat a_m, \hat c_m), & \text{with probability } \rho(\tau), \\
(a_m, c_m), & \text{with probability } 1 - \rho(\tau).
\end{cases}
\]

\section{Sampling from the model}

After training, generation proceeds entirely in sparse-code space before decoding back to pixels.

\subsection{Quantized regime}

Sample autoregressively:
\[
y_t \sim p(y_t \mid y_{<t}),
\]
subject to the structural constraint that atom and coefficient-bin positions alternate appropriately. Then reconstruct the latent tensor $z_q$ and decode:
\[
\hat x = G_\psi(z_q).
\]

\subsection{Real-valued regime}

At each step, sample an atom
\[
a_t \sim p(a_t \mid a_{<t}, c_{<t}),
\]
then predict or sample a coefficient
\[
\hat c_t \approx \mathbb{E}[c_t \mid a_t, a_{<t}, c_{<t}].
\]

The resulting pair sequence is mapped back to the stage-1 sparse latent and decoded to image space.

\section{Conceptual interpretation}

The model imposes three layers of structure.

\subsection{Convolutional structure}

The encoder-decoder pair extracts a multiscale latent representation adapted to images.

\subsection{Sparse linear structure}

The bottleneck forces local latent content to lie near a union of low-dimensional sparse subspaces generated by the dictionary. For fixed support $S$, the latent lives in
\[
\mathrm{span}(D_S).
\]

Across all supports of size at most $s$, the bottleneck defines a union-of-subspaces model:
\[
\mathcal{M}
=
\bigcup_{|S| \le s} \mathrm{span}(D_S).
\]

\subsection{Autoregressive structure}

Stage 2 learns a probability law over the discrete or hybrid coordinates that index points on this sparse manifold. The full generative model is therefore:
\begin{enumerate}
  \item sample a sparse code sequence from the transformer;
  \item map that sequence to a sparse latent point on $\mathcal{M}$;
  \item decode the latent point back to image space.
\end{enumerate}

\section{Why the quantized and real-valued variants differ}

\subsection{Quantized variant}

Everything is turned into symbols. This gives a clean categorical model:
\[
p(y_{1:T}) = \prod_t p(y_t \mid y_{<t}).
\]

Advantages:
\begin{itemize}
  \item simpler likelihood;
  \item simpler sampling;
  \item no scalar regression instability.
\end{itemize}

Cost:
\begin{itemize}
  \item coefficient values are discretized;
  \item reconstruction fidelity may be limited by binning resolution.
\end{itemize}

\subsection{Real-valued variant}

Only atom identities are discretized, while amplitudes remain continuous.

Advantages:
\begin{itemize}
  \item potentially better fidelity because coefficients need not be quantized;
  \item more faithful use of sparse coding geometry.
\end{itemize}

Cost:
\begin{itemize}
  \item stage 2 becomes a hybrid discrete-continuous model;
  \item training requires normalization, auxiliary losses, and careful conditioning.
\end{itemize}

\section{Main approximation built into the model}

The deepest approximation in the whole system is the treatment of sparse coding during stage 1. The exact sparse-coding map
\[
z_e \mapsto \alpha^*(z_e; D)
\]
is non-smooth because support selection changes discontinuously. The model therefore uses a surrogate training procedure:
\begin{itemize}
  \item solve for sparse codes with hard OMP;
  \item reconstruct the latent from those codes;
  \item treat the discrete support choice as fixed during backpropagation.
\end{itemize}

\section{Bottom line}

\texttt{proto.py} implements a sparse-latent generative model whose key object is a learned dictionary-based latent manifold. Stage 1 learns
\[
x \mapsto z_e \mapsto z_q \mapsto \hat x
\]
with $z_q$ constrained to be a sparse reconstruction from learned atoms.

Stage 2 learns a prior over the coordinates of that sparse representation: fully discrete coordinates in the quantized regime, or discrete atom identities plus continuous coefficients in the real-valued regime.

The compact summary is
\[
\text{autoencoder backbone}
+
\text{sparse dictionary manifold}
+
\text{autoregressive prior over sparse coordinates}.
\]

\end{document}
